\documentclass[../libro]{subfiles}

\begin{document}

\ifSubfilesClassLoaded{\mainmatter\chapter{行列式}\clearpage}{}

\section{线性方程组}

在接下来的若干节里, 我想用行列式讨论%
线性方程组的解.
这是行列式的一个应用.
或许, 您可以在这些讨论里体会到
``行列式是一个工具''
``行列式是方阵的一个属性''
的意思.
本节, 我们学习一些基本的概念.

设 \(a_1\), \(a_2\), \(\dots\), \(a_n\) 是 \(n\)~个常数,
\(c\) 是常数,
且 \(x_1\), \(x_2\), \(\dots\), \(x_n\) 是 \(n\)~个未知数.
我们说, 形如
\begin{align*}
    a_1 x_1 + a_2 x_2 + \dots + a_n x_n + c
\end{align*}
的式是一个 \emph{\(n\)~元 \({\leq} 1\)~次式}.
若 \(a_1\), \(a_2\), \(\dots\), \(a_n\)
中有一个数不是 \(0\),
我们说, 这个式是一个 \emph{\(n\)~元 \(1\)~次式}.
若 \(a_1\), \(a_2\), \(\dots\), \(a_n\)
全为 \(0\),
那么, 这个式就是一个常数.

形如
\begin{align*}
    a_1 x_1 + a_2 x_2 + \dots + a_n x_n + c = 0
\end{align*}
的方程是一个 \emph{\(n\)~元 \({\leq} 1\)~次方程}.
不过, 习惯地, 我们移 \(0\)~次项 \(c\) 到等式的右侧;
也就是, 形如
\begin{align*}
    a_1 x_1 + a_2 x_2 + \dots + a_n x_n = b
\end{align*}
的方程也是一个 \(n\)~元 \({\leq} 1\)~次方程,
其中, \(b\) 就是常数 \(-c\).
若我们不想强调未知数数~\(n\),
我们说, 这是一个 \emph{\({\leq} 1\)~次方程}.
我们也可说, 这是一个\emph{线性方程}.

\emph{由 \(m\)~个 \(n\)~元 \({\leq} 1\)~次方程作成的方程组},
是形如
\begin{equation*}
    \left\{
    \begin{aligned}
        a_{1,1} x_1 + a_{1,2} x_2 + \dots + a_{1,n} x_n & = b_1, \\
        a_{2,1} x_1 + a_{2,2} x_2 + \dots + a_{2,n} x_n & = b_2, \\
                                                        & \dots, \\
        a_{m,1} x_1 + a_{m,2} x_2 + \dots + a_{m,n} x_n & = b_m
    \end{aligned}
    \right.
\end{equation*}
的方程组,
其中,
\(a_{1,1}\), \(a_{1,2}\), \(\dots\), \(a_{1,n}\),
\(a_{2,1}\), \(a_{2,2}\), \(\dots\), \(a_{2,n}\),
\(\dots\),
\(a_{m,1}\), \(a_{m,2}\), \(\dots\), \(a_{m,n}\),
\(b_1\), \(b_2\), \(\dots\), \(b_m\)
是事先指定的 \(mn + m\) 个数,
且 \(x_1\), \(x_2\), \(\dots\), \(x_n\)
都是未知数.
若我们不想强调方程数~\(m\),
我们可说, 这是一个 \emph{\(n\)~元 \({\leq} 1\)~次方程组}.
若我们不想强调未知数数~\(n\),
我们可说, 这是一个%
\emph{由 \(m\)~个 \({\leq} 1\)~次方程作成的方程组}.
既然 \({\leq} 1\)~次方程的另一个名字是线性方程,
我们也可说, 这是一个\emph{线性方程组}.

若数 \(c_1\), \(c_2\), \(\dots\), \(c_n\) 适合
\begin{align*}
    a_{1,1} c_1 + a_{1,2} c_2 + \dots + a_{1,n} c_n & = b_1, \\
    a_{2,1} c_1 + a_{2,2} c_2 + \dots + a_{2,n} c_n & = b_2, \\
                                                    & \dots, \\
    a_{m,1} c_1 + a_{m,2} c_2 + \dots + a_{m,n} c_n & = b_m,
\end{align*}
我们说, \((c_1, c_2, \dots, c_n)\)
是此方程组的一个\emph{解}.
有时, 我们也说, 形如
\(x_1 = c_1\),
\(x_2 = c_2\),
\(\dots\),
\(x_n = c_n\)
的 \(n\)~个等式 (的联合)
是此方程组的一个解.

若 \(c_1 = c_2 = \dots = c_n = 0\),
且 \((c_1, c_2, \dots, c_n)\) 是此方程组的一个解,
我们说, 这是此方程组的\emph{零解}.
若 \(c_1\), \(c_2\), \(\dots\), \(c_n\) 不全为 \(0\),
且 \((c_1, c_2, \dots, c_n)\) 是此方程组的一个解,
我们说, 这是此方程组的一个\emph{非零解}.

\begin{example}
    假定有若干只鸡与若干只兔被关在某处.
    假定每一只鸡有 \(1\)~个头与 \(2\)~只腿;
    假定每一只兔有 \(1\)~个头与 \(4\)~只腿.
    假定, 我们知道,
    这些鸡与兔一共有 \(35\)~个头与 \(94\)~只腿.
    我们能由此算出鸡与兔的数目吗?

    我们代数地思考此事.
    设有 \(x_1\)~只鸡与 \(x_2\)~只兔.
    那么,
    这些鸡与兔%
    一共有 \(x_1 + x_2\)~个头与 \(2x_1 + 4x_2\)~只腿.
    注意, 这二个式都是 \(2\)~元 \({\leq} 1\)~次式.
    我们可列出%
    由 \(2\)~个 \(2\)~元 \({\leq} 1\)~次方程作成的方程组
    \begin{equation*}
        \left\{
        \begin{aligned}
            x_1 + x_2   & = 35, \\
            2x_1 + 4x_2 & = 94.
        \end{aligned}
        \right.
    \end{equation*}
    接下来的问题, 就是解这个方程组.
    不过,
    % 不要急;
    此例的目的是使您了解概念.
    我想在后面讲如何解这个方程组.

    可以验证, \((23, 12)\) 是此方程组的一个解:
    因为 \(23 + 12 = 35\),
    且 \(2 \cdot 23 + 4 \cdot 12 = 94\).
    因为 \(23\), \(12\) 里有一个数不是 \(0\),
    故 \((23, 12)\) 是此方程组的一个非零解.
    我们也可说,
    \(x_1 = 23\),
    \(x_2 = 12\)
    是此方程组的一个 (非零) 解.

    不过, \((12, 23)\) 不是此方程组的一个解:
    虽然 \(12 + 23 = 35\),
    但 \(2 \cdot 12 + 4 \cdot 23 = 116 \neq 94\).
\end{example}

\begin{example}
    考虑%
    由 \(2\)~个 \(3\)~元 \({\leq} 1\)~次方程作成的方程组
    \begin{equation*}
        \left\{
        \begin{aligned}
            x_1 + 2x_2 + 3x_3  & = 0, \\
            4x_1 + 5x_2 + 6x_3 & = 0.
        \end{aligned}
        \right.
    \end{equation*}
    可以验证, 对每一个数~\(k\),
    \((k, -2k, k)\) 是此方程组的一个解:
    \begin{align*}
        k + 2(-2k) + 3k  & = k - 4k + 3k = 0,   \\
        4k + 5(-2k) + 6k & = 4k - 10k + 6k = 0.
    \end{align*}
    当 \(k = 0\) 时, 这是零解;
    当 \(k \neq 0\) 时, 因为
    \(k\), \(-2k\), \(k\) 里有一个数 \(k\) 不是 \(0\),
    故这是一个非零解.

    我们也可说,
    \(x_1 = k\),
    \(x_2 = -2k\),
    \(x_3 = k\)
    是此方程组的一个解.
\end{example}

利用阵的积, 我们可简单地写一个线性方程组.
设
\begin{align*}
    A =
    \begin{bmatrix}
        a_{1,1} & a_{1,2} & \cdots & a_{1,n} \\
        a_{2,1} & a_{2,2} & \cdots & a_{2,n} \\
        \vdots  & \vdots  & {}     & \vdots  \\
        a_{m,1} & a_{m,2} & \cdots & a_{m,n} \\
    \end{bmatrix},
    \quad
    B =
    \begin{bmatrix}
        b_1 \\ b_2 \\ \vdots \\ b_m
    \end{bmatrix},
    \quad
    X =
    \begin{bmatrix}
        x_1 \\ x_2 \\ \vdots \\ x_n
    \end{bmatrix}.
\end{align*}
则, 对不超过 \(m\) 的正整数~\(i\),
\begin{align*}
    [B]_{i,1}
    = {} &
    b_i
    \\
    = {} &
    a_{i,1} x_1 + a_{i,2} x_2 + \dots + a_{i,n} x_n
    \\
    = {} &
    [A]_{i,1} [X]_{1,1} + [A]_{i,2} [X]_{2,1}
    + \dots + [A]_{i,n} [X]_{n,1}
    \\
    = {} &
    [AX]_{i,1}.
\end{align*}
注意, \(AX\) 的尺寸也是 \(m \times 1\), 故
\begin{align*}
    AX = B.
\end{align*}
所以, 我们也说形如 \(AX = B\) 的阵等式
(\(X\) 的元是未知数, 且 \(X\) 是恰有 \(1\)~列的阵)
是一个线性方程组.
相应地, 若 \(n \times 1\)~阵 \(C\)
(其的元跟 \(X\)~的元相比, 自然都是已知数)
适合 \(AC = B\),
我们说, \(C\) 是此方程组的一个解;
若 \(C = 0\)
(也就是, \(C\) 的每一个元都是 \(0\))
是此方程组的一个解,
我们说, \(C\) 是此方程组的零解;
若不等于 \(0\) 的 \(C\)
(也就是, \(C\) 有一个元不是 \(0\))
是此方程组的一个解,
我们说, \(C\) 是此方程组的一个非零解.

\begin{example}
    我们可改写
    \begin{equation*}
        \left\{
        \begin{aligned}
            x_1 + x_2   & = 35, \\
            2x_1 + 4x_2 & = 94.
        \end{aligned}
        \right.
    \end{equation*}
    为
    \begin{align*}
        \begin{bmatrix}
            1 & 1 \\
            2 & 4 \\
        \end{bmatrix}
        \begin{bmatrix}
            x_1 \\
            x_2 \\
        \end{bmatrix}
        =
        \begin{bmatrix}
            35 \\
            94 \\
        \end{bmatrix}.
    \end{align*}
    可以验证,
    \(
    C = \begin{bmatrix}
        23 \\
        12 \\
    \end{bmatrix}
    \)
    是此方程组的一个非零解:
    \(C \neq 0\), 且
    \begin{align*}
        \begin{bmatrix}
            1 & 1 \\
            2 & 4 \\
        \end{bmatrix}
        \begin{bmatrix}
            23 \\
            12 \\
        \end{bmatrix}
        =
        \begin{bmatrix}
            1 \cdot 23 + 1 \cdot 12 \\
            2 \cdot 23 + 4 \cdot 12 \\
        \end{bmatrix}
        =
        \begin{bmatrix}
            35 \\
            94 \\
        \end{bmatrix}.
    \end{align*}

    不过,
    \(
    D = \begin{bmatrix}
        12 \\
        23 \\
    \end{bmatrix}
    \)
    不是此方程组的一个解:
    \begin{align*}
        \begin{bmatrix}
            1 & 1 \\
            2 & 4 \\
        \end{bmatrix}
        \begin{bmatrix}
            12 \\
            23 \\
        \end{bmatrix}
        =
        \begin{bmatrix}
            1 \cdot 12 + 1 \cdot 23 \\
            2 \cdot 12 + 4 \cdot 23 \\
        \end{bmatrix}
        =
        \begin{bmatrix}
            35  \\
            116 \\
        \end{bmatrix}
        \neq
        \begin{bmatrix}
            35 \\
            94 \\
        \end{bmatrix}.
    \end{align*}
\end{example}

\begin{example}
    我们可改写
    \begin{equation*}
        \left\{
        \begin{aligned}
            x_1 + 2x_2 + 3x_3  & = 0, \\
            4x_1 + 5x_2 + 6x_3 & = 0.
        \end{aligned}
        \right.
    \end{equation*}
    为
    \begin{align*}
        \begin{bmatrix}
            1 & 2 & 3 \\
            4 & 5 & 6 \\
        \end{bmatrix}
        \begin{bmatrix}
            x_1 \\
            x_2 \\
            x_3 \\
        \end{bmatrix}
        =
        \begin{bmatrix}
            0 \\
            0 \\
        \end{bmatrix}.
    \end{align*}
    可以验证, 对每一个数~\(k\),
    \(
    C = k\begin{bmatrix}
        1  \\
        -2 \\
        1  \\
    \end{bmatrix}
    \)
    是此方程组的一个解:
    \begin{align*}
        \begin{bmatrix}
            1 & 2 & 3 \\
            4 & 5 & 6 \\
        \end{bmatrix}
        \left(
        k\begin{bmatrix}
                 1  \\
                 -2 \\
                 1  \\
             \end{bmatrix}
        \right)
        = {} &
        k
        \left(
        \begin{bmatrix}
                1 & 2 & 3 \\
                4 & 5 & 6 \\
            \end{bmatrix}
        \begin{bmatrix}
                1  \\
                -2 \\
                1  \\
            \end{bmatrix}
        \right)
        \\
        = {} &
        k
        \begin{bmatrix}
            1 \cdot 1 + 2 \cdot (-2) + 3 \cdot 1 \\
            4 \cdot 1 + 5 \cdot (-2) + 6 \cdot 1 \\
        \end{bmatrix}
        \\
        = {} &
        k
        \begin{bmatrix}
            0 \\
            0 \\
        \end{bmatrix}
        =
        \begin{bmatrix}
            0 \\
            0 \\
        \end{bmatrix}.
    \end{align*}
    当 \(k = 0\) 时, 它是零解;
    当 \(k \neq 0\) 时, 它是一个非零解.
\end{example}

最后, 有一件小事值得提.
前面, 我们写一个%
由 \(m\)~个 \(n\)~元 \({\leq} 1\)~次方程作成的方程组%
为阵等式.
反过来,
设 \(A\) 是 \(m \times n\)~阵,
\(B\) 是 \(m \times 1\)~阵,
且 \(X\) 是未知的 \(n \times 1\)~阵.
那么, 我们也可还原形如 \(AX = B\) 的阵等式为%
由 \(m\)~个 \(n\)~元 \({\leq} 1\)~次方程作成的方程组.
比如, 我们可写
\begin{align*}
    \begin{bmatrix}
        2 & 3  & -5 & 7  & 0   \\
        0 & 1  & 11 & 13 & -17 \\
        4 & -6 & 8  & -9 & 12  \\
    \end{bmatrix}
    \begin{bmatrix}
        x_1 \\
        x_2 \\
        x_3 \\
        x_4 \\
        x_5 \\
    \end{bmatrix}
    =
    \begin{bmatrix}
        7 \\
        8 \\
        9 \\
    \end{bmatrix}
\end{align*}
为
\begin{equation*}
    \left\{
    \begin{aligned}
        2x_1 + 3x_2 - 5x_3 + 7x_4         & = 7, \\
        x_2 + 11x_3 + 13x_4 - 17x_5       & = 8, \\
        4x_1 - 6x_2 + 8x_3 - 9x_4 + 12x_5 & = 9.
    \end{aligned}
    \right.
\end{equation*}
(注意, 我们未写系数为 \(0\) 的项.)

\section{\texorpdfstring{由 \(n\)~个 \(n\)~元
      \({\leq} 1\)~次方程作成的方程组 (1)}%
  {由 n 个 n 元 ≤1 次方程作成的方程组 (1)}}

前面, 我们学习了线性方程组的一些基本的概念.
我们知道, 可用阵等式简单地写一个线性方程组
(或者, 形如 \(AX = B\) 的阵等式就是一个线性方程组).
本节, 我们讨论%
由 \(n\)~个 \(n\)~元 \({\leq} 1\)~次方程作成的方程组.
用阵等式表示这类方程组,
就是 \(AX = B\),
其中, \(A\) 是 \(n\)~级阵,
\(B\) 是 \(n \times 1\)~阵,
且 \(X\) 是未知的 \(n \times 1\)~阵.
\(A\) 是方阵, 故它有行列式.
\(AX = B\) 的解跟 \(A\)~的行列式是否有关系?
此事的回答是 ``是''.
我们讨论一个特别的情形.
这会是有用的.

\begin{theorem}[Cramer \pinjino{kērèimer} 公式, 1]
    设 \(A\) 是 \(n\)~级阵 (\(n \geq 1\)).
    设 \(B\) 是 \(n \times 1\)~阵.
    设 \(X\) 是未知的 \(n \times 1\)~阵.
    若 \(\det {(A)} \neq 0\),
    则线性方程组 \(AX = B\) 有唯一的解
    \begin{align*}
        X = (\det {(A)})^{-1} \operatorname{adj} {(A)}\,B.
    \end{align*}
\end{theorem}

在论证此事前, 我想用一个简单的例助您理解它.

\begin{example}
    设 \(a\), \(b\) 是常数.
    一个 \(1\)~元 \({\leq} 1\)~次方程
    \(ax = b\)
    当然也是一个线性方程组:
    这是方程数为 \(1\) 时的特别情形.
    我们在中学就知道, 当 \(a \neq 0\) 时,
    \(ax = b\) 有唯一的解 \(x = a^{-1} b\).
    一方面, \(x = a^{-1} b\) 是一个解:
    \begin{align*}
        a (a^{-1} b)
        = (a\, a^{-1}) b
        = 1b
        = b.
    \end{align*}
    另一方面, 若数~\(y\) 也适合 \(ay = b\), 则
    \begin{align*}
        y
        = 1y
        = (a^{-1} a)y
        = a^{-1} (ay)
        = a^{-1} b.
    \end{align*}

    我们说, 此事是 Cramer 公式的一个特例.
    我们可写 \(ax = b\)
    为阵等式 \([a]\, [x] = [b]\).
    (注意, \([a]\), \([x]\), \([b]\) 都是 \(1\)~级阵.)
    Cramer 公式说, 若 \(\det {[a]} \neq 0\),
    则 \([a]\, [x] = [b]\)
    有唯一的解
    (注意, \(1\)~级阵的古伴是 \([1]\))
    \begin{align*}
        [x] = (\det {[a]})^{-1}
        \operatorname{adj} {([a])}\,[b]
        = a^{-1}\, [1]\,[b]
        = a^{-1}\, [b]
            = [ a^{-1} b ].
    \end{align*}
    这跟我们已知的结论是一样的.
\end{example}

\begin{proof}
    我们先验证
    \(C = (\det {(A)})^{-1} \operatorname{adj} {(A)}\,B\)
    适合 \(AC = B\):
    \begin{align*}
        AC
        = {} & A\,
        ((\det {(A)})^{-1} \operatorname{adj} {(A)}\, B)
        \\
        = {} & (\det {(A)})^{-1} (A\operatorname{adj} {(A)}\, B)
        \\
        = {} & (\det {(A)})^{-1} (\det {(A)}\, I_n B)
        \\
        = {} & ((\det {(A)})^{-1} \det {(A)})\, (I_n B)
        \\
        = {} & B.
    \end{align*}

    我们再证 \(AX = B\) 至多有一个解.
    设 \(n \times 1\)~阵 \(D\) 也适合 \(AD = B\).
    则
    \begin{align*}
        D
        % = {} & 1 I_n D
        % \\
        = {} & ((\det {(A)})^{-1} \det {(A)})\, (I_n D)
        \\
        = {} & (\det {(A)})^{-1} (\det {(A)}\, I_n) D
        \\
        = {} & (\det {(A)})^{-1} (\operatorname{adj} {(A)}\, A) D
        \\
        = {} & (\det {(A)})^{-1} \operatorname{adj} {(A)}\, (AD)
        \\
        = {} & (\det {(A)})^{-1} \operatorname{adj} {(A)}\,B
        \\
        = {} & C.
        \qedhere
    \end{align*}
\end{proof}

\begin{example}
    考虑线性方程组 \(AX = B\),
    其中,
    \begin{align*}
        A = \begin{bmatrix}
                1 & 1 \\
                2 & 4 \\
            \end{bmatrix},
        \quad
        B =
        \begin{bmatrix}
            35 \\
            94 \\
        \end{bmatrix},
        \quad
        X =
        \begin{bmatrix}
            x_1 \\
            x_2 \\
        \end{bmatrix}.
    \end{align*}
    不难算出
    \begin{align*}
        \det {(A)} = 1 \cdot 4 - 2 \cdot 1 = 2 \neq 0,
    \end{align*}
    所以, 此方程组有唯一的解
    \begin{align*}
        X
        = {} &
        (\det {(A)})^{-1} \operatorname{adj} {(A)}\,B
        \\
        = {} &
        \frac{1}{2}
        \begin{bmatrix}
            4  & -1 \\
            -2 & 1  \\
        \end{bmatrix}
        \begin{bmatrix}
            35 \\
            94 \\
        \end{bmatrix}
        =
        \frac{1}{2}
        \begin{bmatrix}
            4 \cdot 35 - 1 \cdot 94  \\
            -2 \cdot 35 + 1 \cdot 94 \\
        \end{bmatrix}
        \\
        = {} &
        \frac{1}{2}
        \begin{bmatrix}
            46 \\
            24 \\
        \end{bmatrix}
        =
        \begin{bmatrix}
            23 \\
            12 \\
        \end{bmatrix}.
    \end{align*}
\end{example}

我们可进一步地改写 Cramer 公式.
我们先具体地算出
\(\operatorname{adj} {(A)}\,B\)~的每一个元:
\begin{align*}
    [\operatorname{adj} {(A)}\,B]_{i,1}
    = {} &
    \sum_{\ell = 1}^{n}
    {[\operatorname{adj} {(A)}]_{i,\ell} [B]_{\ell,1}}
    \\
    = {} &
    \sum_{\ell = 1}^{n}
    {(-1)^{\ell+i} \det {(A(\ell|i))}\, [B]_{\ell,1}}.
\end{align*}
设 \(A \{i, B\}\) 是%
以 \(B\) 代阵~\(A\) 的列~\(i\) 后得到的阵.
则
\begin{align*}
    [A \{i, B\}]_{\ell,k}
    = \begin{cases}
          [A]_{\ell,k}, & k \neq i; \\
          [B]_{\ell,1}, & k = i.
      \end{cases}
\end{align*}
由此可见,
\(A(\ell|i) = (A \{i, B\})(\ell|i)\)
(\(\ell = 1\), \(2\), \(\dots\), \(n\)).
所以
\begin{align*}
         &
    \sum_{\ell = 1}^{n}
    {(-1)^{\ell+i} \det {(A(\ell|i))}\, [B]_{\ell,1}}
    \\
    = {} &
    \sum_{\ell = 1}^{n}
    {(-1)^{\ell+i} \det {((A \{i, B\})(\ell|i))}\,
    [A \{i, B\}]_{\ell,i}}
    \\
    = {} &
    \det {(A \{i, B\})}.
\end{align*}
从而
\begin{align*}
    [X]_{i,1}
    = {} &
    [(\det {(A)})^{-1} \operatorname{adj} {(A)}\,B]_{i,1}
    \\
    = {} &
    (\det {(A)})^{-1}\, [\operatorname{adj} {(A)}\,B]_{i,1}
    \\
    = {} &
    (\det {(A)})^{-1} \det {(A \{i, B\})}
    \\
    = {} &
    \frac{\det {(A \{i, B\})}}{\det {(A)}}.
\end{align*}
由此, 我们得到 Cramer 公式的另一个形式:

\begin{theorem}[Cramer 公式, 2]
    设线性方程组
    \begin{equation*}
        \left\{
        \begin{aligned}
            a_{1,1} x_1 + a_{1,2} x_2 + \dots + a_{1,n} x_n & = b_1, \\
            a_{2,1} x_1 + a_{2,2} x_2 + \dots + a_{2,n} x_n & = b_2, \\
                                                            & \dots, \\
            a_{n,1} x_1 + a_{n,2} x_2 + \dots + a_{n,n} x_n & = b_n.
        \end{aligned}
        \right.
    \end{equation*}
    (注意, 方程数等于未知数数.)
    记
    \begin{align*}
        A =
        \begin{bmatrix}
            a_{1,1} & a_{1,2} & \cdots & a_{1,n} \\
            a_{2,1} & a_{2,2} & \cdots & a_{2,n} \\
            \vdots  & \vdots  & {}     & \vdots  \\
            a_{n,1} & a_{n,2} & \cdots & a_{n,n} \\
        \end{bmatrix},
        \quad
        B =
        \begin{bmatrix}
            b_1 \\ b_2 \\ \vdots \\ b_n
        \end{bmatrix}.
    \end{align*}
    若 \(\det {(A)} \neq 0\),
    则此线性方程组有唯一的解
    \begin{align*}
        x_i = \frac{\det {(A \{i, B\})}}{\det {(A)}},
        \quad
        \text{\(i = 1\), \(2\), \(\dots\), \(n\)},
    \end{align*}
    其中, \(A \{i, B\}\) 是%
    以 \(B\) 代阵~\(A\) 的列~\(i\) 后得到的阵.
\end{theorem}

\begin{example}
    考虑由 \(2\)~个 \(2\)~元 \({\leq} 1\)~次方程作成的方程组
    \begin{equation*}
        \left\{
        \begin{aligned}
            a_{1,1} x_1 + a_{1,2} x_2 & = b_1, \\
            a_{2,1} x_1 + a_{2,2} x_2 & = b_2.
        \end{aligned}
        \right.
    \end{equation*}
    根据 Cramer 公式, 若
    \(\det {
        \begin{bmatrix}
            a_{1,1} & a_{1,2} \\
            a_{2,1} & a_{2,2} \\
        \end{bmatrix}
    }
    = a_{1,1} a_{2,2} - a_{2,1} a_{1,2} \neq 0\),
    则此方程组有唯一的解
    \begin{align*}
        x_1
        = \frac{\det {
                \begin{bmatrix}
                    b_{1} & a_{1,2} \\
                    b_{2} & a_{2,2} \\
                \end{bmatrix}
            }}{\det {
                \begin{bmatrix}
                    a_{1,1} & a_{1,2} \\
                    a_{2,1} & a_{2,2} \\
                \end{bmatrix}
            }}
        = \frac{b_1 a_{2,2} - b_2 a_{1,2}}
        {a_{1,1} a_{2,2} - a_{2,1} a_{1,2}},
        \\
        x_2
        = \frac{\det {
                \begin{bmatrix}
                    a_{1,1} & b_{1} \\
                    a_{2,1} & b_{2} \\
                \end{bmatrix}
            }}{\det {
                \begin{bmatrix}
                    a_{1,1} & a_{1,2} \\
                    a_{2,1} & a_{2,2} \\
                \end{bmatrix}
            }}
        = \frac{a_{1,1} b_2 - a_{2,1} b_1}
        {a_{1,1} a_{2,2} - a_{2,1} a_{1,2}}.
    \end{align*}
\end{example}

\begin{example}
    考虑线性方程组
    \begin{equation*}
        \left\{
        \begin{aligned}
            x_1 + x_2   & = 35, \\
            2x_1 + 4x_2 & = 94.
        \end{aligned}
        \right.
    \end{equation*}
    因为 \(\det {
        \begin{bmatrix}
            1 & 1 \\
            2 & 4 \\
        \end{bmatrix}
    }
    = 1 \cdot 4 - 2 \cdot 1 = 2 \neq 0\),
    故,
    % 由上个例,
    此方程组有唯一的解
    \begin{align*}
        x_1
        = \frac{\det {
                \begin{bmatrix}
                    35 & 1 \\
                    94 & 4 \\
                \end{bmatrix}
            }}{\det {
                \begin{bmatrix}
                    1 & 1 \\
                    2 & 4 \\
                \end{bmatrix}
            }}
        = \frac{35 \cdot 4 - 94 \cdot 1}{2} = 23,
        \\
        x_2
        = \frac{\det {
                \begin{bmatrix}
                    1 & 35 \\
                    2 & 94 \\
                \end{bmatrix}
            }}{\det {
                \begin{bmatrix}
                    1 & 1 \\
                    2 & 4 \\
                \end{bmatrix}
            }}
        = \frac{1 \cdot 94 - 2 \cdot 35}{2} = 12.
    \end{align*}
\end{example}

一般地, Cramer 公式只是一个理论的公式.
这是因为, 一般地, 计算行列式是复杂的事
(或许, 您还记得,
\(3\)~级阵的行列式的具体的公式含 \(6\)~项,
且 \(4\)~级阵的行列式的具体的公式含 \(24\)~项).
所以, 我们一般用别的方法
(如代入消元法、加减消元法)
解线性方程组.

还是以
\begin{equation*}
    \left\{
    \begin{aligned}
        x_1 + x_2   & = 35, \\
        2x_1 + 4x_2 & = 94
    \end{aligned}
    \right.
\end{equation*}
为例.
由方程~\(1\), 有 \(x_1 = 35 - x_2\).
代入它到方程~\(2\), 有
\(2(35 - x_2) + 4x_2 = 94\),
即 \(2x_2 = 24\).
由此可知 \(x_2 = 12\).
代入它到 \(x_1 = 35 - x_2\),
有 \(x_1 = 23\).
最后, 经验证, \((23, 12)\) 的确是此方程组的一个解.

\end{document}

This material tells
what a system of linear equations is,
and gives Cramer's formula.
The two sections are not marked as "optional" ones,
but one can actually safely skip them
if one just wants to learn determinants.
This material,
along with other two materials about systems of linear equations,
is just a application of determinants.
